% !PRINTED 363
\weberpage{363}
%<<CONT
setzen. Hierin kann keine der drei Constanten $\alpha_1, \beta_2, \gamma_3$ verschwinden. Denn wenn z. B. $\alpha_1 = 0$ ist, so schneiden sich nach (9) die Linien $x_2, x_3, y_1$ auf der Curve $f$, und dieser Schnittpunkt müsste ein singulärer Punkt sein.

Nach der Gleichung (9) ergeben sich aber die Coordinaten der Berührungspunkte von $x_1, x_2, x_3$ aus den folgenden drei Paaren von Gleichungen, von denen jedes Paar zwei Berührungspunkte giebt:
\[
\begin{gathered}
x_1 = 0,\quad x_2(\beta_2 x_2 + \beta_3 x_3) - x_3(\gamma_2 x_2 + \gamma_3 x_3) = 0\\
x_2 = 0,\quad x_3(\gamma_3 x_3 + \gamma_1 x_1) - x_1(\alpha_3 x_3 + \alpha_1 x_1) = 0\\
x_3 = 0,\quad x_1(\alpha_1 x_1 + \alpha_2 x_2) - x_2(\beta_1 x_1 + \beta_2 x_2) = 0.
\end{gathered}
\]
Sollen nun diese sechs Punkte auf einem Kegelschnitte $\varphi = 0$ liegen, so muss $\varphi$, von einem constanten Factor $h$ abgesehen, für $x_1 = 0$ in den linken Theil der zweiten Gleichung des ersten Paares $x_2(\beta_2 x_2 + \beta_3 x_3) - x_3(\gamma_2 x_2 + \gamma_3 x_3)$ übergehen. Bezeichnen wir also mit $a_1, a_2, a_3$ die Co\"efficienten von $x_1^2, x_2^2, x_3^2$ in $\varphi$, so muss
\weq{10}{
\begin{gathered}
a_2 = h_1 \beta_2,\quad a_3 = -h_1 \gamma_3\\
a_3 = h_2 \gamma_3,\quad a_1 = -h_2 \alpha_1\\
a_1 = h_3 \alpha_1,\quad a_2 = -h_3 \beta_2
\end{gathered}
}
sein, und hierin sind $h_1, h_2, h_3$ drei Constanten.

Aus (10) folgt aber:
\[
h_2 = -h_3,\quad h_3 = -h_1,\quad h_1 = -h_2,
\]
und dies wäre nur möglich, wenn $h_1 = h_2 = h_3 = 0$ wäre.
Dies ist aber nur dann der Fall, wenn $a_1, a_2, a_3$ verschwinden, wenn also der Kegelschnitt $\varphi$ durch die Schnittpunkte der drei Geraden $x_1, x_2, x_3$ geht. Er soll aber durch die Berührungspunkte dieser drei Doppeltangenten gehen, die sicher von ihren drei Schnittpunkten verschieden sind. Also ist unsere Annahme als unmöglich nachgewiesen und der Satz 3. bewiesen.

\begin{center}
§. 97.

\gesperrt{Complexpaare und Complextripel.}
\end{center}

Die Sätze des vorigen Paragraphen geben ein vorzügliches Hülfsmittel, um die mannigfachen geometrischen und algebraischen Beziehungen der Doppeltangenten zu erforschen und dar-
%CONT>>
